\section{Challenges \& Future Directions}
\label{sec:conclusion}

Having surveyed the state of event camera vision across perception, reconstruction, understanding, and applications, we now discuss the key open challenges and the most promising directions for future research.


\subsection{Open Challenges}
\label{sec:conclusion_challenges}

\paragraph{Data and Benchmarks.}
The scarcity of large-scale, richly annotated event data remains the single most fundamental bottleneck. The largest event classification dataset (N-ImageNet~\cite{kim2021nimagenet}, 1.78M samples) is two orders of magnitude smaller than ImageNet and four orders smaller than LAION-5B. Event-language datasets are even scarcer: the largest (EventMind, 500K+ instructions) is far from the scale needed to train competitive MLLMs. This data gap limits both supervised training and self-supervised pre-training at scale. Moreover, event data is stored in fragmented formats (\texttt{.aedat}, \texttt{.raw}, \texttt{.h5}, \texttt{.npy}, rosbag) with no community-wide standard, creating unnecessary friction. A unified, large-scale, multi-task event benchmark --- analogous to COCO or KITTI for frame-based vision --- would significantly accelerate progress.

\paragraph{Representation and Architecture.}
Despite a decade of research, there is no consensus on the optimal event representation (\cref{sec:pre_repr}). Voxel grids trade temporal resolution for CNN compatibility; raw event graphs preserve full information but are computationally expensive; spike tensors require specialized hardware. The choice of representation is often made ad hoc, with limited systematic comparison across representations and tasks. On the architecture side, Transformers and state space models show promise but have not yet been systematically compared to recurrent networks and SNNs under controlled conditions, particularly regarding the accuracy--latency--energy trade-off that is critical for real-time deployment.

\paragraph{Sim-to-Real Gap.}
Simulated events generated by ESIM~\cite{rebecq2018esim} or v2e~\cite{hu2021v2e} differ from real sensor output due to simplified noise models. Real event camera noise is signal-dependent, spatially non-uniform, and varies across sensor types. DVS-Voltmeter~\cite{lin2022dvsvoltmeter} introduced physics-based stochastic modeling but is not yet widely adopted. Closing this gap is essential for methods that rely on simulation for training (\eg, synthetic events from video datasets via rpg\_vid2e~\cite{gehrig2020vid2e}).

\paragraph{Evaluation Standardization.}
Different papers use different dataset splits, event representations, accumulation windows, and preprocessing pipelines, making fair comparison difficult. Even on the same dataset (\eg, Gen1 for detection, MVSEC for flow), reported numbers may not be directly comparable. The field would benefit from standardized evaluation protocols with fixed preprocessing, splits, and submission-based leaderboards.


\subsection{Emerging Opportunities}
\label{sec:conclusion_opportunities}

\paragraph{Event-Native Foundation Models.}
All current methods for event understanding adapt models pre-trained on RGB data (\eg, CLIP, SAM, Depth Anything). While effective, this approach inherently cannot capture event-specific temporal structure --- the microsecond-scale dynamics, polarity patterns, and motion-dependent sparsity that make events unique. Training a large-scale foundation model \textit{from scratch} on event data could unlock representations that RGB priors cannot provide. The key barrier is data scale: this requires 10--100$\times$ more event data than currently available, potentially achievable through aggressive use of event simulators and text-guided event generation~\cite{wu2025controlevents,barchid2024texttoevents}.

\paragraph{Long-Range Temporal Reasoning.}
Current event MLLMs (EventGPT~\cite{zhou2025eventgpt}, EventVL~\cite{li2025eventvl}) process short event clips (5--1,000 temporal bins). Real-world event streams are continuous and potentially infinite. Extending understanding to long-range temporal reasoning --- \eg, summarizing minutes of event data, tracking causal chains of actions, or predicting future events --- requires fundamentally new architectures that can maintain context over millions of events. State space models (\cref{sec:paradigm_ssm}) and retrieval-augmented approaches may provide pathways.

\paragraph{Unified Multi-Task Models.}
No existing model handles detection, segmentation, depth estimation, captioning, and grounding from events within a single framework. The success of unified models in RGB vision (\eg, GPT-4o, Gemini) suggests that a ``foundation model for events'' --- accepting event streams and producing any combination of structured outputs and natural language --- is both desirable and, given the progress reviewed in this survey, increasingly feasible.

\paragraph{Embodied Intelligence.}
Event cameras' low latency and HDR capability make them ideal for embodied AI --- robots that must perceive and act in the physical world under challenging conditions. E-VLA~\cite{miao2026evla} demonstrated event-augmented robotic manipulation in complete darkness. However, embodied event-based vision remains vastly underexplored: multi-modal fusion with tactile, audio, and proprioceptive signals; long-horizon planning from event observations; and sim-to-real transfer for event-based policies are all open problems.

\paragraph{Real-Time Edge Deployment.}
A persistent gap exists between research algorithms (which typically require GPU computation) and practical deployment (which demands low power and low latency on edge devices). Neuromorphic processors such as Intel Loihi~\cite{davies2018loihi} and SynSense Speck~\cite{liu2023speck} offer a path forward, but most state-of-the-art methods cannot run on these platforms. Algorithm-hardware co-design --- where architectures are designed from the outset for neuromorphic execution --- is needed to close this gap.

\paragraph{Consumer Electronics.}
Despite their low power consumption, event cameras have not yet appeared in smartphones or consumer cameras. Key barriers include manufacturing cost, lack of killer consumer applications, and the need for processing pipelines optimized for mobile SoCs. Event-based computational photography (\eg, HDR capture, super-slow-motion), always-on gesture recognition, and privacy-preserving home monitoring represent potential consumer use cases that could drive mass-market adoption.

\paragraph{Event-Guided Video Generation.}
The success of diffusion-based reconstruction methods (REVDM~\cite{chen2025revdm}, EGVD~\cite{wang2025egvd}) suggests a broader opportunity: using event cameras as \textit{controllers} for general video generation models. Events could guide temporal dynamics, motion trajectories, and scene changes in generated videos with physical plausibility, opening applications in film production, simulation, and content creation.


\subsection{Conclusion}
\label{sec:conclusion_summary}

\cref{tab:survey_comparison} compares the scope of this survey against prior event camera surveys.

\begin{table}[t]
\centering
\caption{Comparison with prior event camera surveys. \cmark: covered; (\cmark): partially covered; \xmark: not covered.}
\label{tab:survey_comparison}
\resizebox{\linewidth}{!}{
\begin{tabular}{lccccccc}
\toprule
\textbf{Survey} & \textbf{Year} & \textbf{\#Papers} & \textbf{Percep.} & \textbf{Recon.} & \textbf{Underst.} & \textbf{Large Models} & \textbf{3DGS} \\
\midrule
Gallego~\etal~\cite{gallego2022survey}    & 2022 & $\sim$200 & \cmark & \cmark & \xmark & \xmark & \xmark \\
Zheng~\etal~\cite{zheng2024deeplearning}  & 2023 & $\sim$150 & \cmark & \cmark & \xmark & (\cmark) & \xmark \\
Xu~\etal~\cite{xu2025recon3d}             & 2025 & $\sim$100 & \xmark & \cmark & \xmark & (\cmark) & \cmark \\
\textbf{This work}                         & \textbf{2026} & \textbf{400+} & \cmark & \cmark & \cmark & \cmark & \cmark \\
\bottomrule
\end{tabular}
}
\end{table}

This survey has provided a comprehensive review of event camera vision in the era of large models, covering over 400 papers across five pillars: benchmarks and datasets, perception, reconstruction, understanding, and applications. We have documented a field undergoing rapid transformation: from niche sensor research to a vibrant ecosystem where event cameras intersect with foundation models, diffusion priors, multimodal LLMs, and 3D Gaussian Splatting.

The key takeaway is that large-scale pre-trained models have expanded what is possible with event cameras, but they should not be conflated with every recent architecture. Open-vocabulary recognition without event-specific labels~\cite{kong2024openess,cho2025ezsr}, reconstruction that adapts video diffusion priors~\cite{chen2025revdm}, and natural language grounding in event streams~\cite{zheng2025talk2event} directly transfer knowledge learned at scale. In parallel, dedicated advances such as pose-free 3D reconstruction~\cite{huang2025inceventgs} arise from event-specific sensing and optimization. Both directions remain important, while fundamental challenges in data scale, evaluation standardization, and edge deployment persist.

We hope this survey, together with its continuously updated companion repository,\footnote{\url{https://github.com/worldbench/awesome-event-camera-vision}} serves as a useful guide for both newcomers and established researchers navigating this exciting and rapidly evolving field.
